In this study, we address the problem of constant injection, withdrawal and idle cycle lengths, and operational parameters currently used in hydrogen storage surrogate models. Deep learning surrogate models have been a growing area of subsurface modelling research in recent years. However, surrogate models require huge amounts of compute capacity to cover a large range of uncertainty which may not be needed. Underground hydrogen storage, in depleted gas fields, is well-placed to use compute efficiently. Due to its cyclical operating model, similar scenarios occur on a yearly basis. Current surrogate models for underground hydrogen storage have been developed using constant cycle lengths and two dimensional data. For future operational use of surrogate models, they must assess a varying range of cycle lengths and operational parameters. This is due to it being thought that hydrogen storage will be operated in the same way as natural gas storage. We show that natural gas storage cycle lengths are variable. We present a highly accurate three dimensional underground hydrogen storage surrogate model, using a U-Net architecture. Hydrogen saturation plume and pressure have maximum errors of 1.6\% and 0.26\% respectively. Predictions of injection rate and hydrogen stored have errors of 5.3\% and 0.649\%. The findings in this study create a framework for future underground hydrogen storage sites to develop highly accurate surrogate models which will enable operational optimisation and integration into global energy systems. 